When a large language model (LLM) is asked a question whose answer it does not know, does it look for a \emph{loophole}---a response that nominally addresses the question while sidestepping its actual demand---rather than admitting ignorance? Prior loophole work has focused on \emph{goal} conflict, where the model knows the answer but its incentives clash with compliance. The analogous \emph{knowledge}-conflict regime, in which the model lacks the relevant facts but still feels pressed to respond, has no targeted benchmark. We introduce \kcl, a paired-prompt benchmark of 20 templates instantiated as topic-matched knowable and unknowable variants. We collect 600 free-form generations and 600 intent multiple-choice probes from five LLMs spanning frontier, large-open, small-open, and long-chain-of-thought reasoning tiers, score them with a four-category rubric judge, and define an outcome as \unsat if it is either a confident wrong answer or a verbalized dodge. Across all five models the unsatisfactory rate is higher on unknowable than on knowable items (mixed-effects GLM odds ratio $3.24$, $95\%$ CI $[1.21, 8.68]$, $p = 0.020$); two reach Bonferroni-significant per-model McNemar effects. The dominant mechanism, however, is \emph{confident hallucination} (27--38\%), not verbalized reinterpretation (dodge rate $\le 5\%$). A complementary \selfaware probe on genuinely unanswerable questions inverts this pattern: models dodge $87$--$97\%$ of the time and almost never refuse. We propose a two-mode taxonomy that distinguishes implicit hallucination loopholes on factual questions from verbal dodge loopholes on underspecified ones, and show that long-chain-of-thought reasoning roughly halves the unsatisfactory rate.
